\documentclass[11pt,a4paper]{article}
\usepackage[margin=2.5cm]{geometry}
\usepackage{enumitem}
\usepackage{booktabs}
\usepackage{hyperref}
\usepackage{titlesec}
\usepackage{fancyhdr}
\usepackage{xcolor}
\usepackage{tabularx}

\hypersetup{
    colorlinks=true,
    linkcolor=blue!60!black,
    urlcolor=blue!60!black
}

\titleformat{\section}{\large\bfseries}{}{0em}{}
\titleformat{\subsection}{\normalsize\bfseries}{}{0em}{}

\pagestyle{fancy}
\fancyhf{}
\rhead{\small\thepage}
\lhead{\small BA International Studies --- Thesis Guidelines}
\renewcommand{\headrulewidth}{0.4pt}

\setlength{\parindent}{0pt}
\setlength{\parskip}{0.8em}

\begin{document}

\begin{center}
{\Large\bfseries BA INTERNATIONAL STUDIES}\\[0.3em]
{\Large\bfseries THESIS GUIDELINES FOR STUDENTS}\\[0.5em]
{\normalsize Leiden University $\cdot$ Faculty of Humanities}\\[1.5em]
\end{center}

\begin{quote}
\small\textit{This document is adapted from the BAIS Thesis Seminars Guidelines and Grading Rubric for student reference. Some content intended for supervisors and second readers has been removed. For the most current information, always consult your supervisor and check Brightspace.}
\end{quote}

\bigskip
\hrule
\bigskip

\section{1. Aims of BA International Studies}

The Bachelor of Arts in International Studies provides students with the tools to investigate globalization and its regional effects from a humanities perspective. Students study these effects through four disciplinary perspectives: culture, history, politics, and economics, combined with in-depth knowledge of one of eight world regions. The humanities perspective places historical and cultural context at the center of the program and links that context to political and economic conditions. Students learn to apply knowledge from the four disciplinary approaches to a geographical area of their choice with the aid of a language native to that area.

\section{2. Objectives of the Bachelor Thesis}

Based on the knowledge and skills acquired, students will be able to:
\begin{enumerate}[label=\alph*)]
    \item work with research techniques that are current in the discipline(s) covered.
    \item understand advanced academic debates.
    \item report on their studies and research in clear written English.
    \item work and write under time-pressure, and deal with deadlines.
\end{enumerate}

The general academic skills covered by these aims are the following:
\begin{enumerate}[label=\alph*)]
    \item collecting and selecting specialized literature using traditional and electronic methods and techniques.
    \item analyzing and evaluating this literature for quality and reliability.
    \item formulating a well-defined research question based on that assessment.
    \item setting up, under supervision, a study of limited size that uses methods and techniques relevant to the discipline.
    \item formulating a reasoned conclusion on that basis.
    \item explaining research findings in writing in a clear and well-argued way.
\end{enumerate}

\section{3. Supervision}

Students working on their bachelor's thesis are supervised by faculty members both in a classroom context (the thesis seminar) and individually. The seminar concentrates on research skills and writing processes. Substantive guidance is provided through individual supervision. The thesis seminar leader is responsible for that individual supervision and acts as the first reader.

Attending the seminar is mandatory. A thesis can be submitted only if it has been written in the context of a thesis seminar.

The thesis should focus on a regional issue and analyze it within a global context. Alternatively, the thesis may examine a regional issue from the perspective of two disciplinary approaches, including at least one of the stated disciplinary perspective(s) of the seminar.

\section{4. Assignments}

Students submit three assignments before submitting the final version of the thesis. These assignments are a prerequisite for thesis submission and provide the primary opportunity for feedback.

\begin{center}
\begin{tabularx}{\textwidth}{lX}
\toprule
\textbf{Assignment} & \textbf{Details} \\
\midrule
Thesis proposal & Research question and plan (approx.\ 1,200--1,500 words) \\
Literature review & Approx.\ 2,500 words \\
Thesis draft & Minimum format to be decided by the supervisor, but including at least a full chapter next to the literature review \\
\bottomrule
\end{tabularx}
\end{center}

All assignments must be discussed with the supervisor individually as soon as possible after submission. The deadline for all assignments is always Friday, at 23:59, in the appropriate week.

Note that due dates cannot be extended unless serious personal circumstances justify the delay. In such a case the supervisor can give a deadline extension of max.\ 5 working days.

\section{5. Knockout Criteria}

A passing grade depends on the following minimum standards. The thesis:

\begin{enumerate}
    \item Contains a clear academic research question.
    \item Is situated within a relevant academic debate.
    \item (a) Places the analysis of the regionally defined topic in a global perspective, or\\
          (b) Analyses the topic from at least two different disciplinary perspectives.
    \item Accounts for the chosen research method(s) and materials.
    \item Is based on the evaluation of a sufficiently large body of independently collected scholarly literature and/or sources (10--20, depending on whether books and/or articles are discussed).
    \item Contains a well-structured and consistent argument.
    \item Is written in correct English.
    \item Produces a scholarly argument and analysis.
    \item Counts 10,000 words ($\pm$10\%), excluding bibliography and notes.
\end{enumerate}

A failing grade for one criterion cannot be compensated with sufficient grades for others.

The referencing style is Chicago. It is up to the supervisor to determine whether that should be the `notes and bibliography' or the `author-date' version.

\section{6. Assessment Criteria}

The thesis is assessed based on the following criteria in the thesis assessment form:

\begin{itemize}
    \item Knowledge and insight (contents, relation to the field)
    \item Application of knowledge and insight (methodology)
    \item Reaching conclusions (interpretation, argumentation, conclusion)
    \item Communication (writing skills, structure)
    \item Learning skills (process)
\end{itemize}

The final grade for the thesis seminar is determined by the thesis grade. To successfully complete this course, the grade for the thesis needs to be a 6.0 or higher. The supervisor is responsible for feedback on all assignments as well as for the assessment of the final thesis. The second reader only assesses the final thesis.

\section{7. Grading Rubric}

\begin{center}
\small
\begin{tabularx}{\textwidth}{lXXX}
\toprule
 & \textbf{Satisfactory (6)} & \textbf{Good (7)} & \textbf{Very good (8)} \\
\midrule
\textbf{1} & & & \\
\textbf{Knowledge} & & & \\
\textbf{\& Insight} & & & \\
\midrule
\multicolumn{4}{p{\textwidth}}{} \\
\bottomrule
\end{tabularx}
\end{center}

\textit{Marks of 9--10 should only be given in truly exceptional cases: i.e., when a student has produced near-publishable work.}

The final grade is derived from an unweighted average of the marks given for the following four criteria:

\subsection{Knowledge and Insight}

\textbf{Satisfactory (6):}
\begin{enumerate}[label=\Alph*.]
    \item Show a general understanding of the relevant literature.
    \item Provide a reasonably clear research question.
    \item Situate the research question in a relevant and reasonably clear theoretical framework.
\end{enumerate}

\textbf{Good (7):}
\begin{enumerate}[label=\Alph*.]
    \item Show a clear and succinct understanding of the relevant literature and demonstrate gaps therein.
    \item Provide a clear and academically topical research question.
    \item Situate the research question in a clear and appropriate theoretical framework.
\end{enumerate}

\textbf{Very good (8):}
\begin{enumerate}[label=\Alph*.]
    \item Demonstrate a full and insightful understanding of the relevant literature, including gaps and connections between schools of academic knowledge.
    \item Provide a clear, verifiable research question on an academically relevant topic.
    \item Situate the research question in an appropriate, well-organized, and properly understood theoretical framework.
\end{enumerate}

\subsection{Application of Knowledge and Insight}

\textbf{Satisfactory (6):}
\begin{enumerate}[label=\Alph*.]
    \item Outline an understandable methodology and describe how data were collected and why.
    \item Describe the main findings coherently on the basis of the data provided.
    \item Organize its body chapters around the data collected in a general way that connects to the research question.
\end{enumerate}

\textbf{Good (7):}
\begin{enumerate}[label=\Alph*.]
    \item Present a clear and sound methodology and justify how data were collected and why.
    \item Describe the main findings of the thesis on the basis of the data and in line with the stated methodology.
    \item Organize the body chapters clearly around the data and the research question.
\end{enumerate}

\textbf{Very good (8):}
\begin{enumerate}[label=\Alph*.]
    \item Provide a clear, academically grounded methodology and explain how the data collection supports the analysis.
    \item Describe the main findings of the thesis in a clear and convincing manner, on the basis of the research question and methodology.
    \item Organize the body chapters in a soundly structured fashion around the data and research question to lead the reader to the conclusions drawn.
\end{enumerate}

\subsection{Reaching Conclusions}

\textbf{Satisfactory (6):}
\begin{enumerate}[label=\Alph*.]
    \item Clearly state its arguments.
    \item Base its arguments on the data presented in the body chapters.
    \item Link its arguments to the literature review and make a case for academic and/or societal relevance.
\end{enumerate}

\textbf{Good (7):}
\begin{enumerate}[label=\Alph*.]
    \item State its arguments clearly and convincingly.
    \item Clearly base its arguments on the data presented in the body chapters and the theoretical framework outlined earlier.
    \item Connect the arguments made to the literature reviewed and argue effectively for academic and/or societal relevance.
\end{enumerate}

\textbf{Very good (8):}
\begin{enumerate}[label=\Alph*.]
    \item State its arguments with clarity and force.
    \item Base its arguments effectively on the data and theoretical framework presented in the thesis.
    \item Convincingly connect the arguments to current scholarly debates and/or broader social contexts.
\end{enumerate}

\subsection{Communication}

\textbf{Satisfactory (6):} Written in reasonably clear academic English and free of endemic grammatical or spelling errors that hinder understanding. The bibliography, citations, and/or footnotes should be formatted correctly with only minor errors.

\textbf{Good (7):} Written in clear academic English and free of notable grammatical or spelling errors. The bibliography, citations, and/or footnotes should be formatted correctly.

\textbf{Very good (8):} Written in clear academic English and free of any serious grammatical or spelling errors. The bibliography and all citations and/or footnotes must be well formatted.

\bigskip

\textbf{NB:} To qualify as `excellent' or `exceptional' for any of the above criteria, the thesis must represent near-publishable content (i.e., material that could, with the right reworking and revision, form part of a published work) and go far above and beyond the descriptions listed above.

\section{8. Thesis Submission}

To submit their thesis, the student needs to:
\begin{enumerate}
    \item Upload the file via a Brightspace assignment, \textbf{or} email it directly to the supervisor, depending on the supervisor's preference.
    \item Email the file to \href{mailto:bathesis@hum.leidenuniv.nl}{bathesis@hum.leidenuniv.nl} on the same day, \textbf{or} include the latter address cc: in the email sent to the supervisor.
\end{enumerate}

Once the thesis has been evaluated with a passing grade, the student needs to:
\begin{enumerate}[resume]
    \item Upload the file in the university's \href{https://studenttheses.universiteitleiden.nl/}{Thesis Repository} (tab `Submit').
\end{enumerate}

\section{9. Late Submission and Resubmission}

\textbf{Late submission:} Students without an extension who fail to hand in their thesis on or before the original deadline, but still within 5 working days, will receive a grade and feedback on their thesis. The grade may be lowered as a consequence of the longer process of completion.

Students who fail to hand in their thesis within 5 working days may still hand in a first version within 10 working days from the original deadline. However, this first version will count as a resubmitted thesis with a consequential lowering of the grade, and in case of a failing grade, there will be no option of handing in a revised version.

\textbf{Resubmission after a failing grade:} A student who scores an insufficient grade (5.0 or lower) is allowed to resubmit a revised version. The deadline for resubmission is 10 working days after receiving the grade and feedback. In case of thesis resubmission, the grade will be lowered as a consequence of the longer process of completion. This lowering should not result in failing an otherwise passable thesis. When establishing the final grade for the resubmitted thesis, an appropriate deduction of the mark will be made for the process. Students who receive a failing grade for their resubmitted thesis have to take another seminar in the next semester.

\section{10. Plagiarism and Academic Integrity}

Students need to be made aware that both plagiarism and the use of AI to generate thesis text are considered fraud. Any suspicion of plagiarism or AI-generated text is to be reported to the Board of Examiners of International Studies immediately.

Please consult the \href{https://www.organisatiegids.universiteitleiden.nl/en/regulations/general/plagiarism}{Regulations on Plagiarism} and the \href{https://www.organisatiegids.universiteitleiden.nl/en/regulations/humanities/guidelines-for-the-use-of-genai-in-assessment}{Guidelines for the Use of GenAI in Assessment} established by the Faculty of Humanities.

\section{11. Thesis Repository}

Students are required to upload the final version of their thesis (as a .pdf file) to the \href{https://studenttheses.universiteitleiden.nl/}{Student Repository} as part of the graduation procedure as soon as they have received a sufficient thesis grade.

Following the university's policy, bachelor's theses are not made publicly available (while master's theses may). For a bachelor's thesis, the following information will be publicly available and visible to search machines: the author, the year, the title, the subtitle, the faculty and the study program, and the supervisor.

Students writing on a politically, culturally, or socially sensitive topic should discuss with their supervisor a safe formulation of the publicly visible data.

\end{document}
